\subsection{Euler contractions and all finite arities for \texorpdfstring{$K(\Z,2)$}{K(Z,2)}}

\begin{theorem}[C27: \status{Proved}; finite bar extensions for a two-sphere parameter]
\label{sa08:C27}
Let $P=K(\Z,2)$, and let $U(P)$ be the positive
bar closure used in Statement~\ref{sa08:C25}.
Let $m\geq2$. Suppose a strict simplicial partial
PD-DGA map $U(P)\to A$ is prescribed through
arity $m-1$, with parameter weight $w$.
Assume that in weight $mw$:
\begin{enumerate}
\item the vertical augmentation from
$\underline{\Z}$ is an isomorphism on $H^0$,
and the vertical cohomology vanishes in
degrees $1,\ldots,2m$;
\item the rows of form degrees $0,\ldots,2m$
are horizontally acyclic.
\end{enumerate}
Then the prescribed map extends through
arity $m$, retaining all its earlier values.
No flatness hypothesis on the target rows
is required for this assertion.

Consequently a map on the original disk
extends through every fixed finite arity
range for which these finite target bounds
hold. If they hold for every $m$ at one
fixed weight $w$, compatible induction
gives a map of the full bar closure.
The eventual axioms alone do not assert
that one such $w$ exists.
\end{theorem}
\begin{proof}
We first establish an integral contraction
statement for a free abelian parameter
module $M$ with a chosen basis.
Let $y_i$ be the associated coefficient
generators of form degree two, and let
$x_i=[y_i]$ be their bar generators.
For each basis index $i$, let $E_i$
be the derivation that counts the
multiplicity of $y_i$ in a basis
multidegree. It is a special PD
derivation: on a divided power of
an even homogeneous element,
the usual factor $n$ is exactly
the identity
$a\gamma_{n-1}(a)=n\gamma_n(a)$.

We construct a special PD derivation
$h_i$ of form degree $-1$ satisfying
\begin{equation}
 dh_i+h_i d=E_i.
 \tag{B-C6}
\end{equation}
It preserves every basis multidegree.
On coefficient generators set
$h_i(y_j)=\delta_{ij}x_i$.
Use the signed PD/exterior free basis
of Statement~\ref{sa08:C24} for the bar
factor and induct on the bar length
of its generators. The differential
of a generator of length $\ell$ is
an expression in generators of
length strictly less than $\ell$,
with coefficient factors of length
zero. Thus $h_i(dg)$ is already
defined when $g$ is reached.

Put $R=E_i(g)-h_i(dg)$. The induction
hypothesis gives $dR=0$. The explicit
bar contraction $\kappa$ in
Statement~\ref{sa08:D029} preserves
multidegree and increases bar length
by one; it satisfies
$d\kappa+\kappa d=1-\varepsilon$.
The element $R$ has positive arity,
so define
\[
 h_i(g)=\kappa(R).
\]
Then $dh_i(g)=R$, establishing (B-C6)
on this generator. Extend $h_i$ as
a special PD derivation on the free
graded PD algebra. Its values may
involve generators of greater bar
length, but this creates no
circularity: only its values on
the generators in $dg$ were used
in the definition.

The commutator of the two odd special
PD derivations $d,h_i$ is an even
special PD derivation. On an even
element, the two cross terms
$da\,h_i(a)$ and $h_i(a)\,da$
cancel. Hence equality on the
free generators proves (B-C6)
on all of $U(M)$. At maximal
bar length the required primitive
has no possible lower form degree;
the closed element $R$ is then
zero by the same explicit bar
contraction identity. There is no
unavailable negative-degree choice.

Let $V_m(M)\subset U_m(M)$ be the
arity-$m$ part generated by lower
arities, including all forced
products and divided powers, and
let $Q_m=U_m/V_m$. Each $h_i$
preserves arity and is a special
PD derivation, so it preserves
this lower-arity subalgebra.
It therefore descends to $Q_m$.
On a component of multidegree
$(n_i)$, (B-C6) is a chain
homotopy from multiplication
by $n_i$ to zero.

Two consequences are crucial.
First, a component containing a
basis variable exactly once is
integrally contractible. Thus any
nonzero vertical homology component
in arity $m$ uses at most
$\lfloor m/2\rfloor$ distinct
basis variables. Second, summing
the Euler identities gives a
homotopy from $m$ times the identity
of $Q_m$ to zero. All its vertical
homology is consequently killed
by $m$.

For $m>1$, the complex $Q_m$ is
supported in form degrees
$m,\ldots,2m-1$. The coefficient
term of form degree $2m$ is a
divided power of arity-one
generators and is already in $V_m$.
Its terms are free, by the
pointwise signed Lyndon basis.
Its degree-$m$ cohomology is the
kernel in a free abelian group.
It is torsionfree and is also
killed by $m$, so it is zero.
Writing $F^q(M)=H^q(Q_m(M))$, we
have proved
\begin{equation}
 F^q=0\quad(q\leq m),\qquad
 \text{each nonzero basis component of }F^q
 \text{ uses at most }\lfloor m/2\rfloor
 \text{ variables}.
 \tag{B-C7}
\end{equation}

We now evaluate this natural homology
functor on $P=K(\Z,2)$. In simplicial
degree $n$ its standard free basis
is indexed by two-element cut sets
of monotone surjections $[n]\to[2]$.
A multidegree component of $F^q(P_n)$
can involve at most
$\lfloor m/2\rfloor$ of these basis
vectors. Their cut sets cover at
most $2\lfloor m/2\rfloor\leq m$
positions. If $n>m$, at least one
cut position is missed. Every
basis vector in that component
then belongs to the image of the
same degeneracy $P_{n-1}\to P_n$.
The full component belongs to
the induced degeneracy image of
$F^q(P)$: restricting to the
split free submodule on those
basis vectors gives the same
multidegree subcomplex and its
same homology.
Therefore
\begin{equation}
 N_nF^q(P)=0\qquad(n>m).
 \tag{B-C8}
\end{equation}
The contractions used to prove
(B-C7) need not be natural in
the chosen basis; the homology
functor and its degeneracy maps
are natural, which is all this
argument requires.

Every abelian group has projective
dimension at most one over $\Z$.
The identity-disk argument used
in Statement~\ref{sa08:C22}
therefore gives
\[
 \operatorname{pd}_{s\mathrm{Ab}}F^q(P)
 \leq m+1.
\]
Every nonzero such object has
$q\geq m+1$ by (B-C7).
It follows that
\[
 \operatorname{Ext}^{q+t}_{s\mathrm{Ab}}
       (F^q(P),\underline{\Z})=0
       \qquad(t>0).
\]
Filter the bounded complex $Q_m(P)$
by its vertical cohomology objects.
The preceding vanishing for every
layer gives
\begin{equation}
 H^{t>0}R\operatorname{Hom}_{s\mathrm{Ab}}
       (Q_m(P),\underline{\Z})=0.
 \tag{B-C9}
\end{equation}
This argument allows nonzero
torsion in many form degrees;
it does not assume a prime-power
concentration formula.

Finally the prescribed lower-arity
map induces a strict map on $V_m$,
by the pointwise free-subalgebra
argument in Statement~\ref{sa08:C25}.
The obstruction to extending it
over $U_m$ is a degree-one
extension class for the cochain
complex $Q_m$. Its terms have
projective normalization, so
extensions by the acyclic target
rows split in each cochain
degree. Thus the obstruction
lies in
$H^1\operatorname{Tot}\operatorname{Hom}
 (Q_m,A^{mw,*})$.

Apply (B-C2) and (B-C3) to
$Q_m[-1]$, supported in degrees
$m+1,\ldots,2m$. The target bounds
in the statement identify this
group with
$H^1R\operatorname{Hom}
 (Q_m,\underline{\Z})$, which
vanishes by (B-C9). An actual
simplicial cochain extension exists.
Every algebra and PD identity
whose output has arity $m$ was
already fixed on $V_m$, so the
extension is the required partial
PD-DGA map. Iterating proves the
finite-range and conditional
infinite-range assertions.
\end{proof}

The theorem does not establish a bar extension
relative to prescribed maps on a parameter
boundary, nor does it control the coefficient
tails required when the eventual bounds grow
with the arity. Its normalized support
estimate specifically uses the two-cut-set
basis of $K(\Z,2)$. These limits remain
separate from the all-finite-arity conclusion.
